\subsection*{$i=4$, first difference at 6: degrees 20 through 23}
\[
 r_6=0,\qquad r_7=-\alpha_6y_5,\qquad
 \alpha_7=(1-s)\alpha_6x_5.
\]
\subsection*{$i=1$, first difference at 5: degrees 12 through 17}
\[
\begin{aligned}
 r_5&=0,&r_6&=-\alpha_5y_2,&\alpha_6&=\alpha_5x_2,\\
 r_7&=-\alpha_5y_3,&\alpha_7&=\frac{s-m}{s}\alpha_5x_3.
\end{aligned}
\]
\subsection*{$i=2$, first difference at 4: degrees 12 through 16}
\[
\begin{aligned}
 r_4&=0,\qquad r_5=-\alpha_4y_3,\qquad
 \alpha_5=\frac{1-m}{1-s}\alpha_4x_3,\\
 r_6&=\alpha_4\left[-b+(s-m)y_3-\frac{s-m}{1-s}x_3y_3-y_4\right].
\end{aligned}
\]
\subsection*{$i=2$, restarted difference at 6: degrees 16 through 19}
\[
 r_6=0,\qquad r_7=-\alpha_6y_3,\qquad
 \alpha_7=(1-m)\alpha_6x_3.
\]
\subsection*{$i=0$, first difference at 2: degrees 4 through 8}
\[
\begin{aligned}
 r_2&=0,\qquad r_3=-\alpha_2y_1,\qquad
 \alpha_3=\frac{1-s}{1-m}\alpha_2x_1,\\
 r_4&=\alpha_2\left[b+(m-s)y_1+\frac{s-m}{1-m}x_1y_1-y_2\right].
\end{aligned}
\]
In the cases $(i,j)=(2,4)$ and $(0,2)$, these formulas leave the higher
coordinates free. If $\alpha_4=0$ in the first case, all difference
coefficients below index 6 vanish; if $\alpha_2=0$ in the second, all
coefficients below index 4 vanish. The argument can therefore restart at
those indices, as in Section~\ref{sec:suff}.
